\cotefile{157-5}
\scanpage{357}\mspage{\deleted{\unsure{3|2}}\unsure{2|3}73}
% slanted
\margin{sous réserve de validité de 4.11.5.}
\uline{Corollaire} 4.11.7. Pour que $M=\unsure{\mathrm{Ind}|\mathrm I\!-\!d}_{\unsure{\mathcal F|\tau}} C$
soit stable par $\underrightarrow{\lim}$ \unsure{lorsque|(resp.} par $\underleftarrow{\lim}$ $J$ \uline{fini},
il \uline{suffit} que $C$ le soit. Et c'est
aussi nécessaire dans le cas envisagé,
lorsqu'on suppose $C$ \unsure{karoubienne}.

4.12. \uline{\deleted{\unsure{Catégories|Remontées|Foncteurs}}\inserted{Foncteurs} accessibles entre catégories
accessibles}. \deleted{\unsure{Sous-catégories}, \inserted{\deleted{\unsure{Par|P}}\unsure{Fam. d'objets|Parties d'objets},} \uline{Accessibilité d'une
catégorie accessible}.}

% slanted
\margin{\unsure{Remonter ceci} à 4.8.}
\uline{Proposition} 4.12.1. Soient $M,N$ deux
catégories accessibles\inserted{ (4.8.)}, et $f:M\unsure{-|\longrightarrow}N$
un foncteur accessible. Alors
il existe \unsure{des|deux} cardinaux \deleted{$\unsure{\pi_1}$}\inserted{infinis $\pi_0,\pi_1$} tels que
\inserted{$\pi_1\geq\pi_0$ et que} pour \deleted{\unsure{tout}} cardinal $\pi\geq\pi_1$, \deleted{\illegible{1}}
on ait
\[
 f(M_\pi)\subset N_{\pi^{\pi_0}}, \tag*{\unsure{(4.12.1.1|(4.12.1.1)}}
\]
et en particulier
\[
 f(M_\pi)\subset N_\pi\qquad\text{si }\pi=\pi^{\pi_0}. \tag*{(4.12.1.2)}
\]
De façon \inserted{plus} précise, on peut prendre
pour $\pi_0$ un cardinal \deleted{\unsure{tel}} utile pour $M$
% archive sheet number: 356
\scanpage{358}\mspage{274}
\deleted{et pour $N$}\inserted{et tel que $f$ soit $\pi_0$-accessible}
(p.ex. le plus petit cardinal \deleted{\unsure{dit utile}}
de \unsure{type|ce type|l'espèce}). Soit alors $\pi_1$ un cardinal tel
qu'on ait $f(M_{\pi_0})\subset N_{\pi_1}$, \inserted{et $\pi_1\geq\deleted{\unsure{\sup}}\bigl(\operatorname{card\,red}M_{\pi_0}\bigr)$, \deleted{\unsure{accessible}}}
\deleted{\unsure{tel que}} \inserted{et enfin, tel que $N_{\pi_1}$ \deleted{soit une cat. \unsure{gén. finie}} soit ss-cat. \inserted{\unsure{str.|stricte}} \unsure{géné-|gén.}.}
\unsure{Alors} $(\pi_0,\pi_1)$ satisfait aux conditions énoncées.

\uline{Démonstration} \hspace{1em} On aura en effet
\[
\begin{aligned}
 M&\xleftarrow{\approx}\mathrm{Ind}_{\pi_0}(M_{\pi_0})\\
 \deleted{N}&\deleted{\xleftarrow{\approx}\mathrm{Ind}_{\pi_0}(N_{\pi_0}),}
\end{aligned}
\]
et $f$ commute\deleted{\unsure{nt|ant}} aux $\unsure{\underrightarrow{\lim}|\underline{\lim}}$ \unsure{grdes} devant
$\pi_0$, sera connu (à isom unique
près) par sa restriction à $M_{\pi_0}$, laquelle
se factorise par $N_{\pi_1}$
\[
 f\mid M_{\pi_0}: M_{\pi_0}\longrightarrow N_{\pi_1}.
\]
Pour prouver (4.12.1.1), \deleted{posant} $\pi'=\pi^{\pi_0}\geq\pi$
et notant que ${\pi'}^{\pi_0}=\pi_0$
\[
 f(M_\pi)\subset f(M_{\pi'})\overset{?}{\subset}N_{\pi'},
\]
on voit qu'il suffit de prouver
(4.12.1.2) (pour $\pi'$). \inserted{Donc \unsure{OPS} $\pi=\pi^{\pi_0}$ et prouvons (4.12.1.2).} Mais comme
$\pi_1\geq\operatorname{Card\,red}M_{\pi_0}$, on sait que si $\pi$ est
tel que $\pi\geq\pi_1$, $\pi^{\pi_0}=\pi$, alors $M_\pi$
est formé des $\underset I{\underrightarrow{\lim}}\,x_i$, les $x_i\in N_{\pi_0}$,
et $I$ \unsure{grd} devant $\unsure{\pi.|\pi_0}$ et $\operatorname{card} I\leq\pi$.
% archive sheet number: 357
\scanpage{359}\mspage{275}
On a alors ($f$ \unsure{étant|ét.} $\pi_0$-accessible)
\[
 f(x)=\underset I{\underrightarrow{\lim}}\,f(x_i),
\]
les $f(x_i)$ sont dans $N_{\pi_1}$, donc dans $N_\pi$,
or $N_\pi$ est stable par $\unsure{\underline{\lim}|\underrightarrow{\lim}}$ de
cardinal $\leq\pi$ (puisque $N$ est engendrée
par $N_{\pi_1}$, formé\unsure{e} d'objets $\pi_1$ accessibles
et a fortiori $\pi$-accessibles, cf
SGA 4 I 9.9. D'où $f(x)\in N_\pi$ \unsure{c.q.f.d.}

\uline{Corollaire} 4.12.2. \inserted{Supposons que $\pi_0$ soit utile pour $N$, et que $\pi_1\geq\operatorname{Card\,red} N_{\pi_0}$.} Soit $\pi$ un cardinal $\geq\pi_1$,
et tel que $\pi^{\pi_0}=\pi$. Alors $f$ \deleted{est isom}
s'identifie, à équivalences de catégories
et à isom \inserted{de foncteurs} près, au foncteur
\[
 \mathrm{Ind}(f_\pi):\mathrm{Ind}_\pi(M_\pi)\unsure{-|\longrightarrow}\mathrm{Ind}_\pi(N_\pi),
\]
où $f_\pi:M_\pi\longrightarrow N_\pi$ est le foncteur
induit par $f$.

% slanted
\margin{\unsure{remonte|remonter} à 4.7.}
\uline{Proposition} 4.12.3. Soient $C$ une catégorie
équivalente à une petite catégorie, $N$
une \inserted{\unsure{$\mathcal U$-|$\mathscr U$-}}catégorie, $\pi_0$ un cardinal infini.
% archive sheet number: 358
\scanpage{360}\mspage{276}
On suppose que $N$ satisfait $L_{\pi_0}$ (stabilité
par $\unsure{\underline{\lim}|\underrightarrow{\lim}}$ gdes devant $\pi_0$). Soit\deleted{\unsure{nt}} $M=\mathrm{Ind}_{\pi_0}C$,
$\underline{\mathrm{Hom}}_{!(\pi_0)}(M,N)$ la sous-catégorie \inserted{str.} pleine
de $\underline{\mathrm{Hom}}(M,N)$ formée des $f:M\unsure{-|\longrightarrow}N$
\deleted{qui \unsure{commutent aux}}\inserted{$\pi_0$-accessibles}. Alors le foncteur
restriction
\[
 \underline{\mathrm{Hom}}_{!\deleted{\unsure{(}}\pi_0\deleted{\unsure{)}}}(M,N)
 \longrightarrow\underline{\mathrm{Hom}}(C,N) \tag*{(4.12.3.1)}
\]
est une équivalence de catégories.
On \unsure{trouve|donne} un foncteur quasi-inverse
comme composé
\[
\begin{aligned}
 \underline{\mathrm{Hom}}(C,N)&\longrightarrow
 (\underline{\mathrm{Hom}}\,\underbrace{\mathrm{Ind}_{\pi_0}(C)}_{M},\mathrm{Ind}_{\pi_0}(N))\\
 &\longrightarrow\underline{\mathrm{Hom}}(M,N),
\end{aligned}\tag*{(4.12.3.2)}
\]
où la deuxième flèche est le foncteur
composition $F\to L\circ F$ avec le foncteur
limite inductive $L:\mathrm{Ind}_{\pi_0}N\unsure{-|\longrightarrow}N$.
Le foncteur composé (4.12.3.2) \deleted{\unsure{qui aura}}
a une image contenue dans $\underline{\mathrm{Hom}}_{!\pi_0}(M,N)$
et induit une équivalence entre \inserted{cette dernière catégorie et}
$\underline{\mathrm{Hom}}(C,N)$. \deleted{\unsure{et celle}}
% archive sheet number: 359

\scanpage{361}\mspage{277}
En d'autres termes, on peut dire que
le foncteur d'induction canonique
\[
 C\unsure{-|\longrightarrow}\mathrm{Ind}_{\pi_0}C=M
\]
est \uline{2-universel} parmi les foncteurs de
$C$ dans des catégories $N$ satisfaisant $L_{\pi_0}$,
celles-ci \deleted{\unsure{flèches}} étant considérées comme
$0$-objets d'une $2$-catégorie où les
$1$-\deleted{objets}\unsure{flèches} sont les foncteurs $N\longrightarrow N'$
$\pi_0$-accessibles, i.e. commu\deleted{\unsure{tent}}\unsure{tants} aux
$\unsure{\underline{\lim}|\underrightarrow{\lim}}$ \unsure{gdes|grdes|grandes} devant $\pi_0$.

Cet énoncé est l'analogue, \inserted{pour $C$ \inserted{petite} quelconque,} de 4.4.16
(page \unsure{101}), valable pour $C$ stable par
$\underline{\lim}$ de cardinal $\leq\pi_0$. Il y a bien
sûr d'autres variantes, en relation avec
la \deleted{stabilité} commutation à d'autres
types de $\underline{\lim}$ ($\underline{\lim}$ filtrantes, sommes,
limites inductives finies etc) envisagés
dans la section 4.11. précédente.

La dém. de 4.12.3 est \unsure{triviale}.
% archive sheet number: 360

\scanpage{362}\mspage{278}
\uline{Corollaire} \deleted{\unsure{3|2}}4.12.4. Soient $M$ et $N$ deux
catégories accessibles, et soit $\pi_0$ un cardinal
infini tel que $M$ soit $\pi_0$-accessible, i.e.
que $\pi_0$ soit un cardinal utile pour $M$.
Soit $\pi_1=\operatorname{Sup}(\pi_0,\operatorname{Card\,red}M_{\pi_0})$. Alors
pour tt cardinal \inserted{$\pi\geq\pi_1$} utile pour $M$
(et en particulier, pour tt $\pi\geq\pi_1$ tel
que $\pi^{\pi_0}=\pi$), la catégorie des foncteurs
$M\unsure{-|\longrightarrow}N$ qui sont $\pi$-accessibles est
elle-\unsure{$=$} accessible.

En effet, on aura $M\approx\mathrm{Ind}_\pi M_\pi$,
d'où par 4.12.3
\[
 \underline{\mathrm{Hom}}_{!\pi}(M,N)\xrightarrow{\approx}\underline{\mathrm{Hom}}(M_\pi,N),
\]
or $\underline{\mathrm{Hom}}(M_\pi,N)$ est accessible par
4.9.5 (\unsure{p|p\textsuperscript{e}|p\textsuperscript{ge}} 220). On peut préciser
d'ailleurs ce point. Supposons $\pi_0$ choisi
tel que $\pi_0$ soit utile non seulemt pour
$M$, mais pour $N$ égalemt, \deleted{et supposons de
\unsure{$=$} $\pi_1\geq\operatorname{Sup}(\operatorname{Card\,red}M_{\pi_0},\operatorname{Card\,red}N_{\pi_0})$.}
\deleted{\unsure{Notons}}\inserted{Rappelons}
que l'on a $\operatorname{Card\,Fl} M_\pi\leq 2^\pi$ (cf. \unsure{4.2.7.1|4.2.7.})
page 173)
% archive sheet number: 361

\scanpage{363}\mspage{27\deleted{\unsure{8}}9}
Posons \deleted{\unsure{donc}}
\[
 \pi'=2^\pi,
\]
on a donc ${\pi'}^{\pi_0}=2^{\pi\pi_0}=2^\pi=\pi'$, et \deleted{\unsure{on a}}
$\operatorname{Card\,red}M_\pi\leq\pi'$. Donc 4.9.5. nous donne
que
\[
 \underline{\mathrm{Hom}}(M_\pi,N)\simeq\mathrm{Ind}_{\pi'}(\mathrm{Hom}(M_{\unsure{\pi}},M_{\pi'})),
\]
d'où

\uline{Corollaire} 4.12.5. Sous les conditions de
4.12.4, si $\pi_0$ est utile pour $M$ \uline{et}
pour $N$, alors pour $\pi\geq\pi_1$, \deleted{\unsure{on aura}}
$\underline{\mathrm{Hom}}_{!\pi}(M,N)$ est \deleted{$\unsure{\mathrm I}$}$\pi'$-accessible, où $\pi'=2^\pi$,
de façon plus précise on a une
\unsure{équivalence} de catégories
\[
 \underline{\mathrm{Hom}}_{!\pi}(M,N)\unsure{\underline{\approx}|\xrightarrow{\approx}}
 \mathrm{Ind}_{\pi'}(\underline{\mathrm{Hom}}(M_\pi,N_{\pi'})).
\]

\uline{Remarque} 4.12.6. Si \deleted{\unsure{cat}} $M,N$ sont accessible\unsure{s},
la catégorie $\underline{\mathrm{Hom}}_{\mathrm{acc}}(M,N)$ des foncteurs
accessibles n'est pas \deleted{\unsure{forcément}} tjrs accessible.
Si $N$ satisfait $L_{\pi_0}$, cette catégorie
satisfait $L_{\pi_0}$, et les $\unsure{\underrightarrow{\lim}|\underline{\lim}}$ gdes devant
$\pi_0$ se calculent argument par argument.
% archive sheet number: 362

\scanpage{364}\mspage{280}
Si on pouvait trouver une sous-catégorie
\inserted{(équiv. à une petite catégorie)}
$\unsure{\mathscr H|\mathcal H}_\pi$ de $\mathscr H=\underline{\mathrm{Hom}}_{\mathrm{acc}}(M,N)$, \inserted{\unsure{pour} un cardinal $\pi\geq\pi_0$,} telle que tt
objet $F$ de $\mathscr H$ s'écrive $\underset{\unsure{\mathbb I|I|J}}{\unsure{\underline{\lim}|\underrightarrow{\lim}}}\,F_i$, les
$F_i$ dans $\mathscr H_\pi$, $\mathbb I$ grd devant $\pi$, alors
\unsure{\{|\deleted{à}} si les objets de $\mathscr H_\pi$ sont \deleted{\unsure{$\pi$-access}}\inserted{des}
foncteurs $\pi'$-accessibles de $M$ dans $N$\unsure{\}}, 
il en serait de $=$ de \unsure{tout|tous} $F$ dans $\mathscr H$.
Donc il existerait un cardinal \inserted{infini} $\pi'$
tel que tout foncteur accessible de
$M$ dans $N$ soit $\pi'$-accessible. Mais
sauf erreur, si \inserted{p.ex.} $M=N=\mathrm{Ens}$, il n'existe
pas de tel cardinal. Je ne veux
pas m'arrêter ici sur cette question,
et faire un contre-exemple.
% archive sheet number: 363

\scanpage{365}\mspage{281}
4.13. \uline{\deleted{\unsure{Parties}}\unsure{Ensembles} accessibles d'objets d'une
catégorie accessible}.

\uline{Définition} 4.13.1. Soient $M$ une catégorie
accessible, $\deleted{M}\unsure{S|\mathscr S}$ une partie de $\operatorname{Ob}M$,
$\underline{\deleted{M}S}$ la sous-catégorie pleine de $M$ définie
par $M$. On dit que $\deleted{M}S$ est une
\uline{partie accessible} de $\operatorname{Ob}M$, ou encore
que $\underline{\deleted{M}S}$ est une \uline{\inserted{ss-}catégorie} ``\uline{accessible
dans $M$}'', si a) $\underline{\deleted{M}S}$ est \unsure{strictemt|strictmt|strictement} pleine,
\unsure{(i.e.|ie} pour \deleted{$x\in M$, et $y\in\operatorname{Ob}M$,} $x,y$
deux objets isomorphes de $M$, $x\unsure{+|\in}\deleted{M}S$
implique $y\in\deleted{M}S$), \deleted{et} b) $\underline{\deleted{M}S}$ est
accessible, c) le foncteur d'inclusion
$\underline{\deleted{M}S}\unsure{-|\longrightarrow}M$ est accessible.

Soit\deleted{\unsure{ent}} $\pi_0$ un cardinal utile pour
$M$, i.e. tel que $M\unsure{\underline{\approx}|\simeq}\mathrm{Ind}_{\pi_0}(M_{\pi_0})$,
soit $\pi_1=\operatorname{Sup}(\pi_0,\operatorname{Card\,red} M_{\pi_0})$. Soit
d'autre part $\pi'_0$ un cardinal tel que
% archive sheet number: 364
\scanpage{366}\mspage{282}
\deleted{\unsure{Il|$\mathrm{Ind}$} \unsure{résulte|suffit} que}
$\underline{\deleted{M}\unsure{S|\mathscr S}}$ (supposée accessible dans $M$) soit
\uline{$\pi'_0$-accessible}, i.e. $\pi'_0$ utile pour $\underline{\deleted{M}S}$, et
de plus tel que $f:\underline{\deleted{M}S}\subset M$ soit $\pi'_0$-accessible.
\deleted{Donc} Soit $\pi'_1$ \inserted{$\geq\pi'_0$, \deleted{\unsure{plus}}} un cardinal tel que
$\deleted{f(}\underline{\deleted{M}S}_{\pi'_0}\deleted{)}\subset M_{\pi'_1}$, $\pi'_1\geq\operatorname{Card\,red}\underline{\deleted{M}S}_{\pi'_0}$.
\deleted{\unsure{et enfin} \illegible{1} et enfin tel que $\pi'_1\geq\pi_0$ \unsure{Il résulte}
Soit enfin $\pi$ un cardinal tel que $\pi\geq\pi'_1$}
\deleted{Notons que, $f$ étant $\pi'_0$-accessible, \inserted{et pl. fid.}
il s'ensuit que \unsure{$S_{\pi}$} pour tt cardinal
$\pi\geq\pi'_0$, on a
\[
 S_\pi\supset S\cap M_\pi
\]}
Il résulte de
4.12.1. que pour tt cardinal
\[
 \pi\geq\pi'_1,\quad\text{tel que }\pi^{\pi'_0}=\pi,\tag*{\unsure{($\alpha$)|(a)}}
\]
on doit avoir
\[
 \underline S_\pi\subset M_\pi\qquad\text{donc }\underline S_\pi\subset M_\pi\cap\underline S
\]
d'ailleurs, comme $\underline S\subset M$ est $\pi'_0$-accessible
et a fortiori $\pi$-accessible, \inserted{et pl. fid.,} il s'ensuit
aussitôt que $\unsure{N|M}_\pi\cap S\subset S_\pi$, d'où\unsure{:}
\[
 \deleted{\unsure{M}}S_\pi=M_\pi\cap S.\unsure{,}
\]
% archive sheet number: 365

\scanpage{367}\mspage{283}
D'ailleurs, les conditions \unsure{($\alpha$)|(a)} impliquent que
$\pi\deleted{\unsure{?}}$ est utile pour $\underline S$, donc
\[
 \underline S\xleftarrow{\approx}\mathrm{Ind}_\pi(S_\pi)\unsure{,|;}
\]

% margin note written sideways
\margin{``facultatif''}
[D'ailleurs, \deleted{\uuline{\unsure{comme}}}\inserted{si} \uline{$\pi'_0\geq\pi_0$}, la relation
$\pi^{\pi'_0}=\pi$ implique a fortiori $\pi^{\pi_0}=\pi$,
et \deleted{\unsure{bien sûr}}\inserted{si} \uline{$\pi'_1\geq\pi_1$}, on voit que
($\alpha$) implique que $\pi$ est utile aussi
pour $M$, i.e.
\[
 M\xleftarrow{\approx}\mathrm{Ind}_\pi(M_\pi).
\]
D'autre part, le foncteur d'inclusion
$\underline S\subset M$ se déduit \inserted{\unsure{essentielt|essentiellement}} de $\underline S_\pi\subset M_\pi$
en passant aux catégories $\mathrm{Ind}_\pi$.]

\margin{$\unsure{\surd}$} Ceci prouve \inserted{entre autres} le ``il faut'' dans la
proposition qui suit\unsure{:}

\uline{Proposition} 4.13.2. Soient $M$ une
cat. \deleted{\unsure{accessible}}\inserted{accessible}, $\pi_0$ un cardinal infini
\deleted{supposons}\inserted{utile pour} $M$, \deleted{$\pi_0$-accessible}. $\unsure{S|\mathscr S}$ une
\deleted{\unsure{sous-caté}} partie de $\operatorname{Ob}M$, \deleted{Pour que}\inserted{$\underline S$ la}
\deleted{$S$ soit accessible}\inserted{sous-catégorie \unsure{pleine} de $M$ correspondante},
% archive sheet number: 366

\scanpage{368}\mspage{284}
a) Pour que $S$ soit accessible, il f. et il
suffit qu'il existe un cardinal \inserted{infini} $\pi$\deleted{\unsure{, tel que} \unsure{($\alpha$)|(a)},
utile pour $M$,} et une partie
$S_{(\pi)}$ de $M_\pi$, \deleted{stable \unsure{dans}} tels que
$S$ soit formé des objets de $M$ qui
sont isomorphes à un objet de la
forme $\underset{\unsure{j|i|v}\in I}{\underrightarrow{\lim}}\,x_i$, avec les $x_i$ dans $S_\pi$,
et $I$ grd devant $\pi$. En d'autres
termes, il f. et s. qu'il existe \deleted{un
$\pi$ \inserted{$\geq\pi_0$} utile pour $M$,}\inserted{un cardinal \unsure{infini} $\pi$,} et une \deleted{\unsure{part.}} ss-catégorie
pleine $\underline S(\pi)$ de $M_\pi$, tels que
$\underline S$ soit l'image essentielle du
foncteur canonique \deleted{\unsure{2}} \unsure{pl fid|pl. fid.|pl fidèle}
\[
 \mathrm{Ind}_\pi(\underline S(\pi))\longrightarrow M
 \bigl(\xleftarrow{\approx}\mathrm{Ind}_\pi(M_\pi)\bigr)
 \quad\inserted{\text{(si }\pi\text{ utile pour }M\text{)}}\tag*{(4.13.2.1)}
\]
prolongeant l'inclusion canonique
$\underline S(\pi)\unsure{\hookrightarrow|\subset} M$ (cf. 4.12.3).

b) \inserted{Si $S$ est accessible,} On peut trouver des cardinaux
$\pi'_0,\pi'_1$, tels que $\pi'_1\geq\pi'_0\geq\pi_0$
\deleted{avec $\pi'_0\geq\pi_0$, $\pi'_1\geq\pi_1=\unsure{\operatorname{Sup}(\pi_0,\operatorname{Card\,red} M_{\pi_0})}$,
$\unsure{\pi'_0,\pi'_1\geq\operatorname{Sup}(\pi_0,\pi_1)}$, $\operatorname{Card\,red}M_{\pi_1}$}
\inserted{(où $\pi_0$ est un cardinal donné utile pour $M$)}
% archive sheet number: 367
\scanpage{369}\mspage{285}
\inserted{et} tels que \deleted{pour} tt cardinal $\pi\geq\pi'_1$ avec
$\pi^{\pi'_0}=\pi$, \deleted{on ait} \unsure{satisfasse} aux conditions
de a), en prenant $S(\pi)=S\cap\operatorname{Ob}M_\pi$,
i.e. $\underline S(\pi)=\underline S\cap M_\pi$. \deleted{Ces catégories
$S(\pi)$ satis\unsure{font} \illegible{1}} Sous ces conditions, \inserted{$\pi$ est utile pour $\underline S$ \uline{et} pour $M$,}
$\underline S(\pi)$ est une ss-catégorie \uline{str.} pleine de $M_\pi$,
stable \inserted{dans $M_\pi$,} par facteurs directs (\unsure{ou ce|ce}
qui revient au $=$, $M_\pi$ étant karoubienne,
une ss-catégorie karoubienne de $M_\pi$).

c) On peut prendre \deleted{\unsure{pour}} pour $\pi'_0$ le plus petit cardinal \inserted{$\geq\pi_0$}
\deleted{\illegible{1}} utile pour $\underline S$ \inserted{et tel que $\underline S\hookrightarrow M$ soit $\pi'_0$-accessible},
\deleted{et enfin $\pi'_1$ le plus petit cardinal $\geq\unsure{\pi'_0,\pi_1},\ \geq\operatorname{Sup}(\pi'_0,\operatorname{Card\,red}\underline S_{\pi'_0})$}
% slanted
\inserted{et $\pi'_1=\operatorname{Sup}(\pi'_0,\unsure{\operatorname{card\,red}\underline{\deleted{M}\unsure{S|\mathscr S}}_{\pi'_0}})$.}

La partie b) \inserted{ainsi que c)} a déjà été prouvée, ce
qui implique le ``il faut'' dans la partie
a). Il reste à prouver le ``il
suffit'' de a). Mais l'hypothèse
implique que
\[
 \underline S\xleftarrow{\approx}\mathrm{Ind}_\pi(S_\pi),
\]
ce qui implique que $\underline S$ est accessible
et $=$ $\pi$-accessible, et elle implique
aussi que $\underline S$ est une sous-catégorie
\uline{str.} pleine de $M$. Il reste à voir
% archive sheet number: 368
\scanpage{370}\mspage{286}
\unsure{Ce|le} foncteur d'inclusion est accessible,
\unsure{en fait} il est $\pi$-accessible, puisque
(4.13.2.1) l'est (cf 4.12.3).

\uline{Définition} 4.13.3. Soient $M$, $S\subset\operatorname{Ob}M$, \inserted{$\underline S\subset M$,} comme
dans 4.12.1, \inserted{$\pi$ un cardinal infini.} On dit que $S$ est
une partie \uline{$\pi$-accessible} de $\operatorname{Ob}M$,
ou que $\underline S$ est une sous-catégorie
$\pi$-accessible dans $M$, si \deleted{$\pi\geq\pi_0$ \unsure{déf}}
\deleted{\unsure{1)|a)} $\pi$ est \unsure{quel|plus} cardinal
plus petit cardinal utile pour $M$,
et si de plus} \deleted{\unsure{b)}} $\unsure{S|C},\pi$ \unsure{satisfont|satisfait} les
conditions de 4.13.3 a), i.e.
$\underline S$ est l'image essentielle de $\mathrm{Ind}_\pi(S(\pi))$,
où $S(\pi)$ est une sous-catégorie pleine
de $M_\pi$.

Notons que, quitte à remplacer
$S(\pi)$ par $S'(\pi)=$ \deleted{sous-catégorie \unsure{st}}\inserted{\unsure{$\leftrightarrows$-plétion|complétion} de} $S(\pi)$
dans $M_\pi$, OPS que $S(\pi)$ est une
sous-catégorie \uline{str. pleine} \uline{et karoubienne}
% archive sheet number: 369
\scanpage{371}\mspage{287}
de $M_\pi$. On sait \deleted{que \unsure{al.} alors que}\inserted{donc que, $S(\pi)$ étant karoubienne,}
$S(\pi)\unsure{-|\longrightarrow}\bigl(\underbrace{\mathrm{Ind}_\pi(S(\pi))}\bigr)_\pi$ est une équivalence
de catégories. Comme $\mathrm{Ind}_\pi(S(\pi))\approx\underline S$,
\deleted{\unsure{al|et}}\unsure{et} que $S(\pi)$ est str. pleine dans
$M_\pi$, donc aussi dans $M$ et a fortiori
dans $\underline S$, que l'on a
\[
 \begin{aligned}
 \underline S(\pi)=\underline S_\pi,\quad&
 \inserted{\text{sous-catégorie de }\underline S\text{ des}}\ \text{objets }\pi\text{-accessibles}\\
 &\text{de }S,
 \end{aligned}
 \tag*{\unsure{(4.13 1)|(4.13.1)}}
\]
Je dis qu'on aura aussi
\[
 \underline S(\pi)=\underline S\cap M_\pi.\tag*{(4.13.1)}
\]
Comme $\subset$ est clair, il reste à voir
que $\underline S\cap M_\pi\subset\underline S(\pi)$, or il est
immédiat que $\underline S\cap M_\pi\subset S_\pi$, et
on conclut par \unsure{(4,13.1)|(4.13.1)}. On \unsure{trouve|trouvera}
ainsi\unsure{:}

\uline{Proposition} 4.13.4. Soit\deleted{\unsure{nt}} $M$ une
catégorie accessible, \deleted{$\pi_0$ le plus petit
cardinal infini utile pour $M$,} $\pi$ un
cardinal \deleted{$\geq\pi_0$}\inserted{infini}. Soit $\Sigma_\pi$ l'ens des
% archive sheet number: 370
\scanpage{372}\mspage{288}
sous-catégories \unsure{strictemt|strictmt|strictement} pleines \inserted{$S(\pi)$} de $M_\pi$
qui sont \deleted{\unsure{stables|sta}}\unsure{karoubiennes}, ou ce qui
revient au $=$, stables \deleted{par} dans $M_\pi$ par
facteurs directs. Soit $\Sigma'_\pi$ l'ens
des sous-\unsure{catégorie|catégories} de $M$ qui sont
$\pi$-accessibles dans $M$. Considérons
les applications
\diagram{\begin{tikzcd}[ampersand replacement=\&]
\text{(4.13.4.1)} \& \Sigma_\pi \arrow[r,shift left=.7ex,"\alpha_\pi"] \& {\Sigma'_\pi} \arrow[l,shift left=.7ex,"\beta_\pi"]
\end{tikzcd}}
définies ainsi : Si $S(\pi)\deleted{\text{\unsure{est}}}\in\Sigma_\pi$,
$\alpha_\pi(S(\pi))$ est l'image essentielle du
foncteur \inserted{\unsure{can} pl fid} $\mathrm{Ind}_\pi(S(\pi))\unsure{\longrightarrow|\text{---}} M$.
Si $\underline S\in\Sigma'_\pi$, $\beta_\pi(\underline S)=\underline S\cap M_\pi$.
Les applications $\alpha_\pi,\beta_\pi$ sont des
bijections, inverses l'une de
l'autre. \deleted{De plus, si $\underline S$ et $S$}

\deleted{\unsure{On n'a pas}} \inserted{On n'a pas} eu à supposer que
$\pi$ soit utile pour $M$, ni que \deleted{\unsure{l'on}}$\pi$\deleted{\unsure{ait}}
\deleted{ait $\pi\geq\pi_0$, où $\pi_0$ est le plus petit
cardinal utile} \inserted{soit \unsure{minoré|majoré} par un cardinal utile pour $M$,}
i.e. qu'on ait $\pi\geq\pi_0$, où
$\pi_0$ est le plus petit cardinal utile pour $M$
% archive sheet number: 371
